% Source: https://artofproblemsolving.com/community/c3h1883859_tricky_geo_problem
% Stable identity: aops-topic-1883859-q01
Triangles $BAD$ and $CAD$ have the same height from $A$, so the ratio of their areas is $\dfrac{BD}{DC}$. We have $\angle CAD=60^\circ-15^\circ=45^\circ$, $\angle ADB=180^\circ-60^\circ-15^\circ=105^\circ$ and $\angle ADC=75^\circ$.

\begin{asy}
usepackage("amsmath");
size(6cm);
defaultpen(fontsize(10pt));
pair A=(0,sqrt(3)),B=(-1,0),C=(1,0);
pair D=(-sqrt(3)*Cos(255)/Sin(255),0);
draw(A--B--C--cycle,linewidth(1));
draw(A--D);
draw(arc(A,0.7,240,255));draw(arc(A,0.3,255,300));
label("$15^\circ$",A+1.0*dir(247.5),fontsize(8pt));
label("$45^\circ$",A+0.5*dir(281),fontsize(8pt));
label("$A$",A,N);label("$B$",B,SW);label("$C$",C,SE);label("$D$",D,S);
label("\mbox{}",(min(currentpicture,user=true)+max(currentpicture,user=true))/2,black); // keeps the text after the figure black when LaTeX typesets the labels inline
\end{asy}

By the law of sines in triangles $ABD$ and $ACD$,
\[
BD=\dfrac{AB\sin15^\circ}{\sin105^\circ},\qquad DC=\dfrac{AC\sin45^\circ}{\sin75^\circ}.
\]
Since $AB=AC$ and $\sin105^\circ=\sin75^\circ$,
\[
\dfrac{BD}{DC}=\dfrac{\sin15^\circ}{\sin45^\circ}.
\]
Now $\sin15^\circ=\sin(45^\circ-30^\circ)=\dfrac{\sqrt2}{2}\cdot\dfrac{\sqrt3}{2}-\dfrac{\sqrt2}{2}\cdot\dfrac{1}{2}=\dfrac{\sqrt6-\sqrt2}{4}$, so
\[
\dfrac{[BAD]}{[CAD]}=\dfrac{\sqrt6-\sqrt2}{4}\cdot\dfrac{2}{\sqrt2}=\boxed{\dfrac{\sqrt3-1}{2}}\approx0.366 .
\]
